\section{Proofs}

This section contains proofs omitted from the main body of this paper and the
lemmas required for these proofs.

\begin{lemma}
  \label{lem:expo-seq}
  The reward for executing $n_I$ improvements and $n_R$ reward actions (in any
  order) is always $\prod_{i < n_I} (1.1 + 0.1 \cdot i)^{e_i}$, where the $e_i
    \in \mathbb{N}$ are exponents, and $n_R = \sum_i e_i$.
\end{lemma}

\begin{proof}
  The universe's transition mechanics imply that a class's reward is a
  product, with the factors being $(1.1 + 0.1 \cdot i)$, one for each
  reward action. $i$ is determined by the number of improvements that have been
  executed before the reward action. $e_{i}$ is simply the number of
  reward actions that happened after improvement $i$, which may be $0$ if
  multiple improvements are made consecutively.
\end{proof}

\begin{lemma}
  \label{lem:breakpoint}
  A strategy that executes all its improvements before its reward actions
  achieves a higher reward than a strategy that executes the same actions in a
  different order. A reward-maximizing strategy must therefore have a
  \emph{breakpoint} $b \in \mathbb{N}$; it only executes improvements before
  that breakpoint and only executes reward actions after that breakpoint.
\end{lemma}

\begin{proof}
  By \cref{lem:expo-seq}, we can rewrite the goal as $(1.1 + 0.1 \cdot
    n_I)^{n_R} \ge \prod_{i < n_I} (1.1 + 0.1 \cdot i)^{e_i}$. The former term
  can be rewritten as $\prod_{i < n_I} (1.1 + 0.1 \cdot n_I)^{e_i}$, which
  proves the goal because $n_I \ge i$.
\end{proof}

\begin{lemma}
  \label{lem:breakpoints-increase}
  As the number of transition steps increases, the breakpoint of
  reward-maximizing strategies also increases monotonically; i.e., to be
  reward-maximizing after more and more transitions, a strategy has to execute
  more and more improvement actions. Formally, for any reward-maximizing
  strategy $s$ at $n$ with breakpoint $b$ (from~\cref{lem:breakpoint}), and any
  reward-maximizing strategy $s'$ at $n'$ with breakpoint $b'$, if $n \le n'$
  then $b \le b'$.
\end{lemma}

\begin{proof}
  By contradiction. Assume $b > b'$. $s$ then has a steeper slope than $s'$
  (because it made $b - b'$ more improvements), and it will have a higher reward
  than $s'$ at $n$ (because it is reward-maximizing at $n$), so it will also
  have a higher reward than $s'$ at $n'$, which contradicts the assumption that
  $s'$ maximizes rewards at $n'$.
\end{proof}

\begin{theorem}
  In the Improvement Game, there exists a horizon $h_n$ for every $n$.
\end{theorem}

\begin{proof}
  Choose any reward-maximizing strategy $s$ with breakpoint $b > n$, which
  maximizes the reward at some point in time $n'$. We assign $h_n$ to be $n' -
    n$.

  Because $b > n$, we know that $s$ executes action $I$ at point $n$. Because
  $s$ is reward-maximizing, we know that this action is also equal to the
  potential-maximizing action $a^*_{h_n}(x,S[x=s],t_n(S[x=s]))$. By
  \cref{lem:breakpoints-increase}, for any $h' > h$, an optimal strategy at $n +
    h'$ must have a breakpoint greater than or equal to $b$, so it must also be
  executing $I$ at $n$.
\end{proof}